\section{Integral spanning by products of initial minors}
\label{sm-section}

The polynomial module in Section~\ref{coef:section} has an additional
integral description needed for branching: its elements are integral
linear combinations of products of initial-column minors. The submitted
proof obtains this description from a two-column argument and
semistandard tableaux.\cite{alpo-uniform-source}
We give a direct proof of the required spanning statement. The only
dimension calculation takes place over characteristic zero; the final
subtraction uses monic polynomials and therefore introduces no
denominators.

\subsection{A two-column leading-term lemma}

Let $d_1,\ldots,d_N$ be nonnegative integers. In the polynomial ring
$\mathbb Q[x_1,\ldots,x_N,y_1,\ldots,y_N]$, let $V$ be the
finite-dimensional space in which the total degree in $x_i,y_i$ is
$d_i$ for every $i$. Put
\[
 E=\sum_{i=1}^N x_i\frac{\partial}{\partial y_i},\qquad
 F=\sum_{i=1}^N y_i\frac{\partial}{\partial x_i},\qquad
 H=\sum_{i=1}^N
 \left(x_i\frac{\partial}{\partial x_i}
             -y_i\frac{\partial}{\partial y_i}\right).
\]
These operators preserve $V$ and satisfy $EF-FE=H$.
For an integer $h$, denote the $h$-eigenspace of $H$ by $V_h$.
Thus a monomial $x^\alpha y^\beta$ belongs to $V_h$ precisely when
\[
 \alpha_i+\beta_i=d_i\quad(1\le i\le N),
 \qquad \sum_i(\alpha_i-\beta_i)=h.
\]
Order monomials lexicographically with the variables in the order
\begin{equation}\label{sm-two-order}
                   x_1>\cdots>x_N>y_1>\cdots>y_N.
\end{equation}
The leading monomial is the greatest monomial with nonzero coefficient
in this order.

\begin{lemma}\label{sm-two-column}
Let $h\ge0$, and let $0\ne f\in V_h$ satisfy $Ef=0$.
If $x^\alpha y^\beta$ is its leading monomial, then
\begin{equation}\label{sm-ballot}
                \sum_{j=1}^i\beta_j
                    \le\sum_{j=1}^{i-1}\alpha_j
                    \qquad(1\le i\le N).
\end{equation}
The same assertion holds over any field of characteristic zero.
\end{lemma}
\begin{proof}
We first compute the dimension of the kernel. Complexify $V$ and make
its monomials orthogonal, with
\[
                \|x^\alpha y^\beta\|^2
                           =\prod_i\alpha_i!\,\beta_i!.
\]
Direct comparison of the coefficients of two monomials gives
$E^*=F$. For example, the contribution involving index $i$ uses
\[
 \beta_i(\alpha_i+1)!\,(\beta_i-1)!
                  =(\alpha_i+1)\alpha_i!\,\beta_i!.
\]
For $v\in V_{h+2}$, the commutator identity consequently gives
\[
            \|Fv\|^2=\|Ev\|^2+(h+2)\|v\|^2.
\]
Since $h+2>0$, the map $F:V_{h+2}\longrightarrow V_h$ is
injective. Its adjoint $E:V_h\longrightarrow V_{h+2}$ is therefore
surjective. The matrices have rational coefficients, so their ranks
over $\mathbb Q$ and $\mathbb C$ agree. It follows that
\begin{equation}\label{sm-kernel-dimension}
       \dim_{\mathbb Q}\ker(E|_{V_h})
                       =\dim_{\mathbb Q}V_h-\dim_{\mathbb Q}V_{h+2}.
\end{equation}

We count the difference on the right without a decomposition theorem
for representations. For $a\ge0$ define the Laurent polynomial
\[
                     \chi_a(T)=T^a+T^{a-2}+\cdots+T^{-a}.
\]
An elementary coefficient calculation gives
\begin{equation}\label{sm-character-product}
              \chi_a(T)\chi_b(T)
                 =\sum_{s=0}^{\min(a,b)}\chi_{a+b-2s}(T).
\end{equation}
Indeed, for $0\le t\le a+b$, the coefficient of $T^{a+b-2t}$
on either side is $1+\min(a,b,t,a+b-t)$; all other exponents have
zero coefficient. On the left this counts pairs $(j,k)$ with
$0\le j\le a$, $0\le k\le b$ and $j+k=t$. On the right it
counts $s$ with $0\le s\le\min(a,b)$ and $s\le t\le a+b-s$.

The dimension of $V_h$ is the coefficient of $T^h$ in
$\prod_i\chi_{d_i}(T)$. Iterating \eqref{sm-character-product},
starting from $h_0=0$, expresses this product as a sum of the
polynomials $\chi_{h_N}(T)$, one for each sequence
\begin{equation}\label{sm-path}
 0\le\beta_i\le\min(d_i,h_{i-1}),\qquad
 \alpha_i=d_i-\beta_i,\qquad
 h_i=h_{i-1}+d_i-2\beta_i
 \quad(1\le i\le N).
\end{equation}
Every $h_i$ is nonnegative. For $h\ge0$, the coefficient of $T^h$
minus the coefficient of $T^{h+2}$ in $\chi_k(T)$ is one if $k=h$
and zero otherwise. Therefore \eqref{sm-kernel-dimension} is the
number of sequences \eqref{sm-path} with $h_N=h$.
Since
\[
                    h_{i-1}=
                         \sum_{j<i}(\alpha_j-\beta_j),
\]
these sequences are exactly the nonnegative exponent arrays with the
prescribed row degrees and total weight $h$ which satisfy
\eqref{sm-ballot}.

For each such array we construct a monic polynomial in the kernel
with leading monomial $x^\alpha y^\beta$. Process the indices
$i=1,\ldots,N$ in order. Before processing $i$, retain a product of
brackets $x_jy_k-x_ky_j$ with $j<k<i$, together with
$h_{i-1}$ unpaired factors $x_j$, all with $j<i$.
Choose $\beta_i$ of the unpaired factors and replace each chosen
$x_j$ by
\[
                            x_jy_i-x_iy_j.
\]
Then adjoin $\alpha_i$ factors $x_i$. The choice is possible because
$\beta_i\le h_{i-1}$; the number of unpaired factors afterwards is
$h_{i-1}-\beta_i+\alpha_i=h_i$. This starts with the empty product
at $i=1$, when \eqref{sm-ballot} says $\beta_1=0$.

Every bracket and every unpaired $x_j$ is killed by $E$, so their
final product is killed by $E$. Its row degrees and weight are the
prescribed ones. Because $j<i$, the leading term of the displayed
bracket in \eqref{sm-two-order} is $x_jy_i$, with coefficient one.
Consequently the final product has leading monomial
$x^\alpha y^\beta$ and leading coefficient one.
Choose one such product for each exponent array satisfying
\eqref{sm-ballot}. Distinct choices of array give distinct leading
monomials, hence linearly independent polynomials. Their number is
the dimension in \eqref{sm-kernel-dimension}; they form a
$\mathbb Q$-basis of the kernel.

In a nonzero linear combination of these polynomials, the greatest
of their leading monomials with nonzero coefficient cannot cancel:
every monomial in the other participating polynomials is smaller.
This proves \eqref{sm-ballot} for the leading monomial of $f$.
Finally, the matrix of $E$ and all the constructed polynomials have
integer coefficients. The rank of that matrix is unchanged under
extension from $\mathbb Q$ to any characteristic-zero field.
The same products are therefore a basis of the corresponding kernel
over that field, and the leading-monomial conclusion is unchanged.
\end{proof}

\subsection{Several columns and integral subtraction}

Fix integers $N\ge r\ge1$ and a partition
$\mu_1\ge\cdots\ge\mu_r\ge0$. Let $X=(x_{ij})$ be an $N$ by $r$
matrix of indeterminates. For a characteristic-zero integral domain
$R$, define $\nabla_\mu(R)$ to be the set of polynomials in $R[X]$
which are homogeneous of degree $\mu_j$ in each column $j$ and
satisfy
\begin{equation}\label{sm-column-invariance}
          f\bigl(X(I+tE_{lj})\bigr)=f(X)
                    \qquad(1\le l<j\le r)
\end{equation}
as polynomial identities in the additional indeterminate $t$.
Here $E_{lj}$ is the elementary matrix, so this substitution adds
$t$ times column $l$ to column $j$.
Use the lexicographic order
\begin{equation}\label{sm-column-order}
 x_{11}>x_{21}>\cdots>x_{N1}>
 x_{12}>\cdots>x_{N2}>\cdots>x_{1r}>\cdots>x_{Nr}.
\end{equation}
For a strictly increasing row sequence
$I=(i_1<\cdots<i_s)$, $1\le s\le r$, put
\[
               \Delta_I(X)=
                 \det(x_{i_a,b})_{1\le a,b\le s}.
\]
This is an initial-column minor of height $s$.

\begin{lemma}\label{sm-leading-interlacing}
If $0\ne f\in\nabla_\mu(R)$ has leading monomial
$\prod_{i,j}x_{ij}^{a_{ij}}$, then for $j<r$ and $1\le i\le N$,
\begin{equation}\label{sm-interlacing}
                   \sum_{b\le i}a_{b,j+1}
                         \le\sum_{b<i}a_{b,j}.
\end{equation}
\end{lemma}
\begin{proof}
Pass to the fraction field of $R$. Differentiating
\eqref{sm-column-invariance} at $t=0$ shows that
\[
                       E_jf=0,\qquad
              E_j=\sum_i x_{ij}\frac{\partial}{\partial x_{i,j+1}}.
\]
Decompose $f$ according to the monomial in all columns other than
$j,j+1$ and according to the degrees
$d_i$ in $x_{ij},x_{i,j+1}$. The operator $E_j$ preserves each
of these finite components, and hence kills each one.
In particular it kills the nonzero component containing the leading
monomial of $f$. After removing its fixed monomial in the other
columns, this component has the two-column degrees $\mu_j,\mu_{j+1}$,
so its weight is $h=\mu_j-\mu_{j+1}\ge0$.
The restriction of \eqref{sm-column-order} to the two columns is
exactly \eqref{sm-two-order}. Its leading monomial is the part of
the leading monomial of $f$ in these two columns. Applying
Lemma~\ref{sm-two-column} gives \eqref{sm-interlacing}.
\end{proof}

\begin{lemma}\label{sm-matching-minors}
Suppose a nonnegative integer array $(a_{ij})$ has column sums
$\mu_j$ and satisfies \eqref{sm-interlacing}. There is a product of
initial-column minors with the same column degrees, leading monomial
$\prod_{i,j}x_{ij}^{a_{ij}}$, and leading coefficient one.
It has exactly $\mu_s-\mu_{s+1}$ factors of height $s$, where
$\mu_{r+1}=0$.
\end{lemma}
\begin{proof}
For each column $j$, form the nondecreasing list
\[
                       w_j(1),\ldots,w_j(\mu_j)
\]
in which the row index $i$ occurs $a_{ij}$ times. If
$1\le t\le\mu_{j+1}$ and $w_{j+1}(t)=i$, there are at least
$t$ entries of the list for column $j+1$ at most $i$.
By \eqref{sm-interlacing}, the list for column $j$ has at least
$t$ entries strictly less than $i$. Thus
\[
                         w_j(t)<w_{j+1}(t).
\]
For each $1\le t\le\mu_1$, put
\[
 s_t=\max\{j:t\le\mu_j\},\qquad
 I_t=(w_1(t),\ldots,w_{s_t}(t)).
\]
The preceding inequalities make $I_t$ a strictly increasing row
sequence. Consider $P=\prod_{t=1}^{\mu_1}\Delta_{I_t}(X)$.
In the determinant of any initial-column minor, its diagonal term
is the greatest term in \eqref{sm-column-order}: column one chooses
the smallest available row index, then column two chooses the
smallest remaining one, and so on. Its coefficient is one.
Leading monomials and coefficients multiply, so $P$ has the claimed
leading monomial and leading coefficient one. A factor has height
$s$ exactly when $\mu_{s+1}<t\le\mu_s$; this proves the last
assertion as well as the column degrees. When $\mu_1=0$, all lists
are empty and $P=1$ has the required properties.
\end{proof}

\begin{theorem}[Integral spanning by initial minors]
\label{sm-initial-minors}
For $N\ge r\ge1$ and a partition $\mu$ as above,
$\nabla_\mu(\mathbb Z)$ is spanned over $\mathbb Z$ by products
of initial-column minors having exactly $\mu_s-\mu_{s+1}$ factors
of height $s$, for $1\le s\le r$ and $\mu_{r+1}=0$.
The same equality of modules holds over any characteristic-zero
integral domain $R$, with the identical integral polynomials as
generators and coefficients in $R$.
\end{theorem}
\begin{proof}
An initial-column minor of height $s$ is unchanged under every
substitution in \eqref{sm-column-invariance}. If $j>s$, column $j$
does not occur. If $j\le s$, both columns $l<j$ occur, and the
substitution is a determinant-preserving column operation.
The indicated product of minors therefore belongs to
$\nabla_\mu(R)$, since its degree in column $j$ is
\[
                     \sum_{s\ge j}(\mu_s-\mu_{s+1})=\mu_j.
\]

Conversely, let $f\in\nabla_\mu(R)$ be nonzero, with leading
coefficient $c\in R$ and leading exponent array $(a_{ij})$.
Lemma~\ref{sm-leading-interlacing} and
Lemma~\ref{sm-matching-minors} supply a product $P$ of the required
minors with that leading monomial and leading coefficient one.
Then $f-cP$ still belongs to $\nabla_\mu(R)$ and, unless it is
zero, has a strictly smaller leading monomial.
There are only finitely many monomials with the fixed column degrees
$\mu_1,\ldots,\mu_r$. Repeating this subtraction must therefore
terminate at zero. It expresses $f$ as a linear combination over
$R$ of the specified integral products. Every subtraction used a
leading coefficient in $R$ and a monic polynomial; no division in
$R$ occurred. Taking $R=\mathbb Z$ proves the integral assertion.
\end{proof}

The case of no columns consists of constants and is represented by
the empty product. Zero parts of $\mu$ likewise cause no exception.
This argument proves exactly the integral spanning input and its
characteristic-zero extension. It does not assert that differentiating
the column identities suffices in positive characteristic: in
characteristic $p$, for example, $y^p$ is killed by
$x\partial_y$ but is not invariant under $y\mapsto y+tx$ as a
polynomial identity. The submitted source states a more general
result for domains of arbitrary characteristic; the present proof
does not replace that additional assertion. All coefficient lattices
to which we apply Theorem~\ref{sm-initial-minors} are in
characteristic zero.


